\subsection{Societal Impact}
\label{sec:societal-impact}

MERIT changes the perception of ``good evidence''. Its namesake was an intentional reference to the concept of meritocracy.
While this may be considered a positive outcome in some cases, it is important to note that it is a trade-off that is inherent in the design of MERIT.

Relying less on CV polish and ATS-friendly formatting can help prevent candidates from being penalised for not having a perfect CV.
However, it may also lead to a shift from keyword stuffing on CVs towards gaming an online portfolio. Candidates could introduce
commit streaks, bloat repositories, and inflate activity metrics to appear more active. Alongside this, candidates who work primarily
in private repositories (such as enterprise, finance and defence) can score lower than juniors with active public GitHub profiles,
despite a clear difference in tenure. This is a trade-off that is inherent in the design of MERIT.

MERIT does not fix hiring on its own. Fair outcomes depend on employers keeping the human veto, not forcing GitHub portfolios as part of
an application process, and monitoring who gets rejected when evidence is missing.
Table~\ref{tab:ethics-deployment-summary} summarises the main risks, what MERIT does about them, and what problems remain.

\begin{table}[htbp]
    \centering
    \caption[Deployment risks and mitigations]{Deployment risks, MERIT mitigations, and residual issues.}
    \label{tab:ethics-deployment-summary}
    \small
    \renewcommand{\arraystretch}{1.15}
    \begin{tabularx}{\textwidth}{@{}%
        >{\raggedright\arraybackslash\hsize=0.95\hsize}X
        >{\raggedright\arraybackslash\hsize=1.15\hsize}X
        >{\raggedright\arraybackslash\hsize=0.9\hsize}X @{}}
        \hline
        \textbf{Risk} & \textbf{What MERIT does} &
            \textbf{What remains} \\
        \hline
        Rejecting people with no human review (Art.~22) &
        Scores support recruiters, shortlist needs a person
            (Sec.~\ref{sec:fairness-computational-rigour}) &
        Illegal if fully automated \\
        \hline
        Bias from names, schools, titles &
        Blind Mode, smaller prestige gap in Study~15 &
        Scores still shown, tier labels kept \\
        \hline
        Users do not trust the system &
        Glass-Box scores and Shapley breakdown &
        Over-trusting scores, detail unused \\
        \hline
        Keyword stuffing on CVs &
        Anti-stuffing, multiple data sources &
        Gaming GitHub activity (Study~06) \\
        \hline
        Fake profiles / stolen GitHub &
        Integrity checks, Study~14 caught fraud &
        Rare smart squatting, name mismatches \\
        \hline
        Penalising private-sector work &
        GitHub optional, CV still used &
        ``Proprietary gap'' (Carl Corp) \\
        \hline
        LinkedIn scraping (prototype only) &
        Apify, at most 30 consenting contacts
            (Sec.~\ref{sec:data-privacy-gdpr}) &
        Need licensed API in production \\
        \hline
    \end{tabularx}
\end{table}
